\documentclass[11pt]{article}
\usepackage{amsmath,amssymb,amsthm}
\usepackage[margin=1.15in]{geometry}
\newtheorem{theorem}{Theorem}
\newtheorem*{slogan}{Slogan}
\newcommand{\Sym}{\mathrm{Sym}}
\title{Multi-Dimensional Harmonic Number Theory,\\ for Dummies}
\author{Samuel Lavery}
\date{August 1, 2026 --- draft}
\begin{document}
\maketitle

\begin{abstract}
This is a plain-language tour of a working method: read the classical objects of
analytic number theory --- $L$-functions, their zeros, their signs, their central
values --- as flat photographs of a higher-dimensional geometric object, and do the
bookkeeping on the object instead of the photograph.  The method's signature is
\emph{exactness}: where the classical field estimates, this one keeps ledgers that
balance to zero on the nose, at a scale set by sixth roots of unity.  No prior
acquaintance with the program is assumed.  Everything here is either proven
(machine-checked or cited), measured (with the instrument stated), or explicitly
open --- and each claim says which.
\end{abstract}

\section{The one-sentence idea}

\begin{slogan}
An $L$-function is a photograph.  The subject of the photograph is a helix carrying
rotating arrows, and arithmetic is easier when you do your accounting on the helix.
\end{slogan}

Take the Riemann zeta function on its critical line.  The classical picture is a
wiggly graph: a one-dimensional signal, with mysterious moments where it vanishes.
Multi-dimensional harmonic number theory (MDHNT) says this signal is the
\emph{readout} of something geometric: a three-dimensional carrier --- a double
helix and its conjugate anti-helix, harmonically scaled, running from the origin
out to infinity --- on which each integer $n$ contributes one \emph{phasor}: an
arrow of length $n^{-1/2}$, spinning at rate $\log n$.  The value of the
$L$-function is what you get when you add up all the arrows.  A zero is a moment
of perfect cancellation: the whole bank of arrows, tip to tail, closing into a
polygon.

None of this is a metaphor bolted on after the fact.  The projection from the
carrier to the classical signal is a bijection \emph{once you keep its ledger}: the
descent $3\mathrm D\to1\mathrm D$ books exactly what it forgets --- a height, a
radius, an angle per datum --- and projection-with-ledger has an exact two-sided
inverse (a machine-checked theorem in the program's Lean corpus).  The photograph
loses information; the photograph \emph{plus the ledger} does not.

Why bother?  Because some things that are hard in the photograph are easy on the
object, and --- this is the working discovery of the program --- some things that
look like \emph{walls} in the photograph are not there at all in three dimensions.

\section{The goal: a Langlands-style toolkit, but executable}

Here is what is actually being built, stated as an engineer would.  The Langlands
program unified number theory by decreeing a common home: lift everything to
automorphic representations, and relationships between arithmetic objects become
morphisms in one category.  It is the great \emph{intermediate representation} of
mathematics --- but it is a correspondence, not a computation surface.  MDHNT's
bet is a common home you can \emph{run things on}: a geometric intermediate
representation with instruments attached.  The toolkit has five stages, and the
reader should hold the pipeline in mind through everything that follows.

\begin{enumerate}
\item \textbf{Lift (harmonize).}  Take an ordinary analytic function --- an
$L$-function, a zeta-like series, in principle anything with a series or product
presentation and growth control --- and homogenize it onto the carrier as a
phasor configuration: a magnitude law, a phase law, a harmonic scale.  Different
inputs become \emph{comparable} because they now live in one format.
\item \textbf{Choose a geometry.}  The carrier is a parameter, not a commitment:
the standard double helix, or non-standard geometries when the problem wants
different structure exposed.  Same configuration, different backend --- like
running one program on different hardware to see different profiles.
\item \textbf{Load multiple rails.}  A carrier holds as many objects as you
want, on parallel rails --- two, ten, a whole family --- and alignment can be
taken over any chosen subset of rails, jointly or pairwise.  This is the
comparison mechanism, and it is strictly stronger than the photograph: the
program has an exhibit (a genuinely non-self-dual $\mathrm{GL}(3)$ object) that
the multi-rail representation carries faithfully while the scalar
one-dimensional readout \emph{provably cannot}.
\item \textbf{Align (register).}  The core primitive.  Bring chosen rails into
correspondence --- clocks synchronized, vanishing points matched --- and read
the \emph{defect ledger} of the alignment, which is integer-valued and exact.
``How are $f$ and $g$ related?''\ becomes: lift both, align, read the ledger.
Shared zeros are aligned cancellation events; lifts and transfers are rail maps;
sign laws are diamond bookkeeping.  Relationship becomes accounting.
\item \textbf{Extract (read out).}  Project the answer back to the classical
one-dimensional chart --- or any other chart --- with the loss ledger making the
extraction exact rather than lossy.  Results leave the toolkit in whatever
language the consumer speaks.
\end{enumerate}

One condition makes the whole pipeline legal, and it is the program's real
signature: \textbf{every stage keeps a ledger, and ledgers compose}.  The lift is
a bijection with its loss ledger; alignment defects are exact; extraction is
invertible-with-ledger.  That is why ``one universal harmonization technique''
can be an honest engineering claim instead of a slogan: modularity is legal
exactly when nothing is silently lost at the joints, and here nothing is ---
by theorem, stage by stage.

The endgame capability, stated plainly: a \emph{calculus of comparisons}.  Given
any collection of objects --- not two, however many you want --- the toolkit
answers ``evaluate these against each other'' by harmonizing all of them and
reading the aligned ledgers, over whichever rails you choose --- a
one-size-fits-all unification of inputs and outputs, as modular as we can make
it.  And admissibility is a spectrum, not a gate.  The lift accepts nearly
anything, including functions with no arithmetic structure at all --- random
inputs are welcome.  What changes with structure is what the ledgers \emph{do}:
structured objects close exactly; a random input produces an unpredictable
ledger --- and that output has value, because it is a \emph{measurement}.  The
random ensemble is the toolkit's calibration class: structure announces itself
as a ledger anomaly against that background, which is how the instrument
carries its own null hypothesis --- the house discipline of controls-that-must-fail,
built into the pipeline itself.  And this claim is pre-registered, not
decorative --- and it has now been \emph{run}, blind, by a cold reader, with
the pre-registration corrected by the data, which is how this method is meant
to work.  The sharp null is a random \emph{completely multiplicative} phase
(Steinhaus: independent uniform phases on the primes, extended
multiplicatively), preserving the multiplicative skeleton and killing only the
arithmetic.  The null must generate \emph{its own events}: load the null fiber
on the double-ended carrier and take the real zeros of the conjugate-pair
completion $2\,\mathrm{Re}\bigl[e^{i\vartheta(t)}F_f(\tfrac12+it)\bigr]$ ---
the reflection symmetry manufactures the functional equation the model lacks.
(An earlier surrogate --- freeze the arithmetic event counts, randomize
phases --- is question-begging: the occupancy sequence is precisely what
carries the signal.)  Verdict, at $128{,}000$ cells and $16$ realizations: six
arithmetic inputs (zeta and five Dirichlet $L$-functions, each on its own
clock) show \emph{zero} registration defects --- not small; zero --- while
the sharp null mis-registers at $\approx0.9\%$ of cells, so its running sum
grows like $N^{1.78}$ by the dichotomy (an unabsorbed defect costs linear
growth; a defect \emph{rate} costs more).  The original $\sqrt N$ prediction
survives exactly where it belonged: the null's pure phase content
random-walks at $N^{0.54}$.  And the measured discriminator is sharper than
anything pre-registered: \emph{the turning-point exclusion zone}.  Arithmetic
banks backtrack as often and as deeply as the null, but keep their turning
points a hard margin ($\approx0.09\pi$) from the clock's level set, where the
null's turning points graze it freely --- an exclusion of probability
$\sim10^{-46}$ under a uniform model.  Registration is arithmetic in the
literal sense: same excursions, different placement against the clock.  (The
weak null --- all phases independent --- fails by \emph{drift}, not random
walk: a cheap control, now for the honest reason.)
The success
criteria are falsifiable, in the house style: the toolkit must re-derive known
correspondences as rail alignments (it has begun to); it must exhibit relations
the one-dimensional chart provably cannot see (it has); and if lifting and
aligning competent inputs consistently produces empty ledgers and no new
relations, that is evidence against the toolkit claim, and will be published as
such.

\section{Chart artifacts, or: the wall that wasn't}

Here is the cheapest example, and the program's founding habit of mind.  Every
textbook says the Dirichlet series for $\zeta$ ``only converges for
$\mathrm{Re}\,s>1$,'' and continuation past that line is treated as a minor
miracle.  On the carrier there is no line and no miracle: each phasor enters
continuously at zero magnitude and grows until finished; the bank exists at every
height; nothing diverges, because nothing is being summed in the order the
photograph suggests.  The convergence abscissa is a fact about the \emph{chart} ---
the coordinate system of the readout --- not about the object.

We call such things \emph{chart artifacts}, and the program maintains a standing
rule: never accept a chart artifact as an obstruction to a statement about the
object.  This rule has teeth.  More than once, a ``gap'' that stalled work for days
was dissolved --- not solved, dissolved --- by checking whether it existed in
carrier coordinates.  (A recent instance: a numerical law ``$|A_n|<1$'' about
certain cell integrals turned out to be the product of a genuine carrier-native law
with a shrinking chart factor.  In carrier coordinates the true law was visible,
sharper, and provable; the chart version was a coincidence of ranges.)

The converse habit matters equally: a fact can be \emph{real but chart-dressed},
and then the chart version is a legitimate readout --- functional equations, local
factors, root numbers all have honest one-dimensional shadows.  The discipline is
to know, for every statement, which side of that line it lives on.

\section{The cast of characters}

\paragraph{The carrier.}  A double helix and its conjugate anti-helix,
harmonically scaled.  ``Harmonically scaled'' is load-bearing: the natural rulers
on the carrier are sixth roots of unity --- $\pi/3$ cells, the lattice
$\mathbb Z[\zeta_6]$ --- and measurements taken at unit scale instead of harmonic
scale produce false nulls.  (This was learned the hard way, twice, and is now a
method law.)

\paragraph{Phasors.}  One per site $n$: magnitude $n^{-1/2}$, phase driven by
$\log n$.  The bank of phasors \emph{is} the arithmetic object; the $L$-function
is its resultant.  Magnitudes enter by a partial-absorption law --- each arrow
grows across its own site and completes there --- so the bank is always a finite,
concrete thing at any height.

\paragraph{The clock.}  A reference rotation $\vartheta(t)$ (in the classical
chart, the Riemann--Siegel theta function) against which everything is registered.
The clock is the carrier's own time.  Measured in clock units $d\vartheta/\pi$,
cells have unit length by definition, and --- a running theme --- integrals that
look transcendental in the chart become exact algebra in the clock.

\paragraph{Ledgers.}  Integer-valued running books.  The classical function
$S(t)$, usually described as the erratic ``argument of zeta,'' is in this
framework a \emph{registration gap}: the running difference between two counts
that ought to agree.  The ledger itself oscillates in sign; what is
non-negative, integer-valued and monotone is the \emph{defect} between the
chart's ledger and the carrier's native one (see the $S(t)$ worked example) ---
a distinction that matters, and that an earlier draft of this sentence blurred
at the cost of one careful reader's detour.  The reframing is a theorem, and it
converts questions about wild analytic objects into questions about whether
books balance.

\section{The object itself, defined}

The previous section named the characters; this one pins them down, because a
reader (human or machine) should be able to check the geometry rather than take
it on faith.  There are three layers, and keeping them distinct dissolves most
confusions.

\paragraph{The geometry (the stage).}  The standard configuration is a
double-ended helix, and here is its actual parametrization (machine-checked,
with its arclength in closed form):
\[
\gamma(k) \;=\; \bigl(r k\cos 2\pi k,\; r k\sin 2\pi k,\; p k\bigr),
\]
Archimedean in the turn parameter $k$ --- radius $rk$ growing linearly, pitch
$p$ --- registered against the analytic ordinate by the \emph{climber}
$k(y)=e^{y}/p$, so that height $z=pk=e^{y}$ and cylindrical radius
$R(y)=(r/p)\,e^{y}$.  Both classical descriptions are therefore true at once,
chart-dependently: the spiral is Archimedean in its own parameter and
exponential in the ordinate chart --- one more instance of the chart lesson.
The carrier's tick marks are arclength in units of $\Delta=\pi/3$ (the
root-of-six system), with the continuous tick index $N(y)=S(k(y))/\Delta$; the
exact mod-6 bucket structure lives on the \emph{integer site spins}
$s_n=n\cdot\pi/3$ (six-periodic, with exact counts per residue class), not on
$N(y)$ --- $N$ carries the configuration factor and nothing arithmetic is read
off it modulo anything.  The partner \emph{anti-helix} is the conjugate,
reflected through the origin; both ends run $0\to+\infty$.  The geometry is
chosen per run and then fixed; nothing downstream rescales it.

\paragraph{The carrier (the ruler).}  The geometry \emph{hosts} a carrier: a
scaled numberline laid along it.  This is where harmonic scaling lives ---
scaling a carrier means multiplying it by a harmonic constant ($\pi/6$,
$\pi/3$, $\pi/2$, \dots), and \emph{warping} it means applying a dynamic
rescaling function; the two can be layered.  One geometry can host many
carriers at once --- these are the \emph{rails}: differently configured
carriers running next to each other in the same space, overlapping without
displacing one another.

\paragraph{The fiber (the function).}  The $L$-function's representation lives
on the carrier, and the phasors live on the fiber.  As the fiber grows from the
origin, the conductor assigns each carrier integer a positive, negative, or
neutral phasor; each is born at magnitude zero and grows to full size, and each
spins \emph{in log} because it is spinning on a rotating carrier --- its phase
is relative to its own start time and to the other phasors.

\paragraph{The descent, and where the ledger entries come from.}  Angle is
always measured relative to the centroid of the configuration.  The projection
to the classical picture is a two-step descent, each step booking exactly what
it drops: $3\mathrm D\to2\mathrm D$ drops the \emph{radius}, by the explicit
Cayley map $w\mapsto(w-i)/(w+i)$ applied to the readout coordinate $w$ --- the
fiber's reading on the real readout line, which the map sends onto the unit
circle --- and
$2\mathrm D\to1\mathrm D$ drops the \emph{angle} (unit circle to the strip,
with the ordinate read as $y=\log z$).  Read that again
with a number theorist's eyes: \emph{the critical strip is the Cayley image of
the unit circle} --- two ledger entries downstream of the object.  That is why
strip phenomena keep turning out to be chart artifacts.

\paragraph{Where the geometry does visible work: the functional equation is a
reflection.}  The deepest identity in the subject --- the functional equation,
$s\leftrightarrow1-s$ --- is, in the flat chart, a minor miracle: it needs theta
inversion, Poisson summation, analytic continuation, and a completed
$\Gamma$-factor, and it arrives as a symmetry with no visible cause.  On the
carrier it has a cause you can point at: \emph{the geometry is double-ended}.
The helix and its conjugate anti-helix are mirror images through the origin,
and the readout cannot tell which end it is reading.  The reflection that swaps
the two ends descends, through the Cayley step, to exactly the strip's
$s\leftrightarrow1-s$, and the arithmetic strand pair
$\mathrm{diag}(\alpha,\alpha^{-1})$ --- one leg on each end --- is the
determinant-one pairing whose trace the chart reads.  Poisson summation belongs
nowhere in this picture: it is the \emph{chart's} device --- the way the
one-dimensional readout reconstructs, from inside the projection, a symmetry
the object simply has.  On the carrier there is nothing to prove; in the chart
there is a theta transformation law to establish, and Poisson summation is how
the chart earns what the geometry was born with.

\paragraph{What the conjugacy endows.}  The right calibration of the thesis is
this: the geometry is not a tool you reach for on every problem --- it is the
system's \emph{endowment}.  The conjugate pairing alone supplies, before any
particular problem is posed: the functional equation (the reflection above);
the \emph{Frobenius similitude} --- the arithmetic strand block
$\mathrm{diag}(\alpha,\alpha^{-1})$ has determinant one \emph{because} the two
legs are conjugate ends, which is what makes the radial squeeze well-posed at
all; the axis-reality of paired readings (Schwarz pairing for free); and the
two-strand completion that dissolves the classical ``homelessness'' of
single-strand objects --- one strand alone has no mirror term and no functional
equation, and a century of orphaned objects were orphans of the chart, not of
the mathematics.  Some results never mention any of this, just as a real
integral evaluated by contour methods never mentions the complex plane in its
statement.  Nobody concludes that $\mathbb C$ is scaffolding.  The geometry may
not be useful for every problem; the system is incomplete without it.  In the compiled corpus
the pair reflection is a theorem, not a picture.  A note on skepticism done
well: a sharp reader's test for geometric content --- ``does $r$ or $p$ survive
into any result?''\ --- detects \emph{shape}-work and is right to apply it; the
functional equation is \emph{symmetry}-work, and symmetry leaves no parameter
fingerprint in a statement.  It leaves the statement itself.

\paragraph{A zero, geometrically.}  A vanishing point is a \emph{focal
cancellation}: the positive and negative phasor channels close against each
other, all pointing to the same \emph{eigenheight} $z$ on the carrier; the
absolute phase energy locates it, the fiber --- itself a harmonic --- flips
sign there, and the \emph{harmonic pencil} pins the exact location.  In
symbols, at a candidate height $Z$, the two finite channel readings of the bank
are
\[
A=\sum_{n<N}\chi(n)\,n^{-(3/2+i\log Z)},\qquad
B=\frac{\pi}{3}\sum_{n<N}\chi(n)\,n^{-(1/2+i\log Z)}
\]
--- the unsigned channel at the absolutely convergent abscissa and the
$\pi/3$-cell--normalized channel at the critical one --- and the pencil is the
matrix family $\begin{pmatrix}A&B\\ \mu A&\lambda B\end{pmatrix}$ over chosen
weights $(\mu,\lambda)$, with determinant exactly $(\lambda-\mu)AB$: degeneracy
at $Z$ \emph{is} the vanishing certificate.  Two finite sums, no continuation;
the $L$-value appears only as final verification.  The
one-dimensional chart reads the event at $y=\log z$: the famous ordinate of a
zero is the logarithm of a carrier height.  Aligning vanishing points
\emph{across rails} is the correspondence signal --- when two differently
loaded carriers cancel at a common height, something more interesting than
either function has happened.

\section{The techniques, as recipes}

\subsection{Registration: make two clocks agree, or price the defect}

The signature move.  Given two bookkeeping systems for the same phenomenon --- a
classical chart count and a carrier-native count --- prove an \emph{exact} law of
the form
\[
\text{(chart count)} \;=\; \text{(native count)} \;+\; \text{(defect)},
\]
where the defect is integer-valued, monotone, and non-negative.  Then the whole
question ``do the systems agree?''\ collapses to ``is the integrated defect small?''
--- and because a single unabsorbed unit of defect costs \emph{linear} growth in
the integral, there is no intermediate regime: sublinear defect is already zero
defect.  (Machine-checked.)  You do not chase individual events; you audit the
books.

\subsection{The clock chart: integrate in carrier time}

Cell-by-cell integrals of a ledger, taken in $dt$, depend on cell lengths and look
analytic.  Taken in $d\vartheta/\pi$ --- carrier time --- they become exact
combinatorics: a cell's ledger entry is a boundary term, plus one term per event, each
an explicit affine function of the event's \emph{clock phase}.  Here is the
entire instrument, in four lines --- this chain \emph{is} the construction:
\[
L(t)=N(t)-\frac{\vartheta(t)}{\pi}-1;\qquad
\text{cell } m=[g_m,g_{m+1}],\ \vartheta(g_m)=m\pi,\ \text{origin } m=0;
\]
\[
r_m=\int_{\text{cell }m} L\,\frac{d\vartheta}{\pi};\qquad
B_N=\sum_{m<N} r_m .
\]
Conventions, so no reader ever reconstructs them again: cells are half-open
$[g_m,g_{m+1})$, a boundary event belongs to the cell on its right, events are
counted with multiplicity, the entering ledger is $s_m=N(g_m)-m-1$, and an
event at $t_i$ in cell $m$ has clock phase $\tau_i=\vartheta(t_i)/\pi-m\in[0,1)$.
The chart's whole content is integer boundary data: the ledger enters cell
$m$ at the integer $s_m$ and exits at the integer $s_{m+1}=s_m+k-1$ --- no
halves anywhere.  Between endpoints, in carrier time, the drift is exactly
linear --- down one unit per cell, up one at each event --- so the integral
$r_m$ is the \emph{average} of that drift, and averaging is where halves are
manufactured, as midpoints of unit descents: an event-free cell books
$s-\tfrac12$, a one-event cell $s+\tfrac12-\tau$, and in general
$r_m=s_m+k-\tfrac12-\sum_i\tau_i$.  The $\tfrac12$ is mean-of-drift ---
$B_1$, the sawtooth's offset --- an averaging artifact that belongs to the
trapezoid, never to the cell.  Carrying the mean back to the boundary
recovers the integers:
$r_m+\sum_i\tau_i+\tfrac12=s_m+k=s_{m+1}+1$ in every cell --- the sharpest
free check any implementation of this instrument can run.  In a recent campaign this
single change of chart turned a fog of numerics into three exact laws in an
afternoon, and the observed pinning of phases traced back, through an exact
identity, to the $n\ge2$ phasors --- the Euler stream itself.

\subsection{Lattice rigidity: machine zero is zero}

The norm form of $\mathbb Z[\zeta_6]$ takes integer values, and its smallest
nonzero value is $1$.  Consequence: any quantity that (a) provably lives on the
lattice and (b) measures below norm $1$ is \emph{exactly zero}.  This is the
program's bridge from computation to proof --- a numerical residual of $10^{-12}$
is not ``small,'' it is a certificate, because the only lattice point in that ball
is the origin.  Kronecker's theorem plays the same role one level up: a lattice
value of modulus one is a sixth root of unity \emph{exactly}, so ``approximately
frozen'' phases are honestly frozen, with an address.

\subsection{Towers: climb with one ruler}

Arithmetic objects come in towers (symmetric powers, tensor powers, grades of
cycles), and the method insists on rank-uniform instruments: one ceiling, one
pairing, one cell decomposition for the whole tower.  Rank-uniformity is not
aesthetics --- it is where limits come from.  Example (one paragraph, complete):
if a strand's radius $\rho$ satisfies $\rho^r\le C$ for every rung $r$ of the
tower --- one fixed ceiling, every rung --- then $\rho\le1$; applied to the strand
and its inverse, $\rho=1$ exactly.  That elementary squeeze, fed with the tower's
automorphy and the classical per-rung bound, is temperedness (the chart calls it
Ramanujan--Petersson), and on the carrier it says: \emph{the strand is a helix,
not a spiral}.

\subsection{The discipline: numerics first, gates before claims, registers always}

Every idea is probed numerically before it is formalized; every instrument
carries validation anchors and is run against controls designed to \emph{fail}
(a control that fails by exhibiting \emph{more} structure than the target is a
design bug --- also learned the hard way).  Every novelty claim passes a
literature gate: the closest papers are read at source, with pin-cites, before
the word ``new'' is typed --- a gate that has both cleared claims and killed
them, sometimes in the same afternoon.  And every statement is published in one
of three registers, never blended: \emph{proven} (machine-checked in Lean at a
fixed axiom footprint, or cited to the classical literature at the point of
use), \emph{measured} (precision-tracked numerics --- corroboration, never
proof), or \emph{program} (stated with its falsifiers, in print, in advance).

\section{Worked example: $S(t)$, the mystery that became a ledger}

The most famous ``erratic'' function in analytic number theory is $S(t)$, the
argument of zeta along the critical line.  It enters the zero-counting formula
\[
N(t)\;=\;1+\frac{\vartheta(t)}{\pi}+S(t),
\]
where $N$ counts zeros up to height $t$ and $\vartheta$ is the smooth clock.
Read that formula the classical way and it says something remarkable that is
rarely said out loud: \emph{the smooth clock does essentially all the work}.
The mysterious part, $S$, is a small correction --- oscillating, mean zero,
unbounded but barely ($O(\log t)$ unconditionally) --- and a century of
literature treats it as the residue of chaos in the zeros.

The carrier reading is different, and it is a theorem, not a mood: $S(t)$ is a
\emph{registration gap} --- the running disagreement between two bookkeeping
systems for the same events, the chart's count and the carrier's native count
against its own clock.  It is not noise about the zeros; it is the ledger of
how the chart's ruler slips against the carrier's.  The native ledger is
\[
S_{\mathrm{native}}(t)\;=\;N_{\mathrm{native}}(t)-1-\frac{\vartheta(t)}{\pi},
\]
a step count minus a smooth ramp --- exactly the shape the clock-chart recipe
of \S5.2 wants --- and the exact registration law
\[
S_{\mathrm{chart}}\;=\;S_{\mathrm{native}}\;+\;D,
\qquad D\ \text{a non-negative, monotone, integer step function,}
\]
is compiled, together with the dichotomy of \S5.1: $D$ has no intermediate
regime --- a sublinear integrated defect is already a zero defect, because one
unabsorbed unit costs linear growth forever.

Now put the recipes to work in the clock chart, cell by cell, and the
``erratic'' function resolves into structure you can hold.  Its true content
is integer boundary data --- the ledger walks from integer to integer, cell
by cell; the halves in the cell means are averaging midpoints of the
interior drift, not properties of the cells.  So an event-free cell's mean
sits at a \emph{floor} --- $-\tfrac12$ plus its boundary defect, with zero
freedom; a one-event cell wears its event's clock phase affinely;
and measured across thousands of cells the books balance absurdly well --- the
running integral of the native ledger stays inside a band of width about
$1.4$ across five thousand cells, the floors' two populations cancel by defect
neutrality, and the event phases are not free: they are slaved, through an
exact identity, to the $n\ge2$ phasors --- the Euler stream itself deciding
where in each cell a zero may sit.

And here, in symbols, is the program's open analytic residual --- the one
statement this framework has reduced everything to, stated so an outsider can
attack it:
\[
\boxed{\ \int_0^{T} S_{\mathrm{native}}(t)\,dt\;=\;o(T).\ }
\]
That is all.  The reduction from the full identification statement to this
single mean bound is machine-checked and lossless; so is the fact that the
current corpus cannot close it without a genuinely new ingredient --- both
facts are part of the record, not embarrassments hidden from it.  The measured
margin is enormous (the integral appears \emph{bounded}, not merely $o(T)$),
the mechanism carrier (the Euler stream's exceedance episodes) is identified,
and the statement contains no zeros, no strip, and no continuation --- it is a
statement about one ledger's mean, in the object's own coordinates.

\section{Worked example: root numbers by carrying digits}

A fresh result, small enough to show whole.  The sign $\pm1$ in an $L$-function's
functional equation decides whether it is forced to vanish at its center.
Classically these signs are computed prime by prime, and for symmetric powers of
an elliptic curve the answer is \emph{periodic} in the power.  Nobody had computed
the sign for the full \emph{tensor} tower $H^1(E)^{\otimes g}$.

On the carrier the archimedean sign of grade $g$ is $i^{A(g)}$, where $A(g)$ is
read off the Hodge diamond, and the diamond sum collapses to a closed form:
\[
A(g)\;=\;2^{\,g-1}+g\binom{g-1}{\lfloor (g-1)/2\rfloor}.
\]
When is the sign $-1$?  Exactly when $A(g)\equiv2\pmod4$; and by Legendre's
formula the $2$-adic valuation of the central binomial $\binom{2m}m$ is the
number of carries when you add $m+m$ in base two (Kummer, 1852).  One carry
exactly means $m$ is a power of two.  Conclusion, machine-checked:
\[
\varepsilon_\infty=-1 \iff g=2^k+1 .
\]
Root numbers, decided by grade-school addition.  The set $\{3,5,9,17,33,\dots\}$
is not periodic --- which is precisely why it could not have come from the
classical per-prime bookkeeping, and why the literature gate found the territory
unclaimed.  Assembling the global sign (multiplicities are ballot numbers;
per-piece signs are the classical periodic ones, cited) yields the payoff: for
every semistable curve, at every Fermat grade $g=2^k+1$, the tensor-tower
$L$-function vanishes at its center to \emph{odd} order --- with computed
multiplicity ($3$ at the Gross--Schoen grade; about $1.9\times10^{18}$ by
$g=65$).  Two different arithmetics --- carry-counting at the archimedean place,
congruence classes at the finite places --- multiplying into one Fermat-shaped
law.

\section{What the method has reached (three registers, as always)}

\emph{Proven.}  The depth dictionary for cycle filtrations; retention on every
finite extension stack; the registration laws for $S(t)$ and the exact
identification criterion; temperedness and Selberg-type conclusions from the
tower's radial limit; Sato--Tate for non-dihedral Maass forms (via the tower,
Jacquet--Shalika, and Serre's criterion --- a case outside the reach of
automorphy-lifting methods, which need a Galois representation the Maass world
does not offer); the freeze laws and their $\mu_6$ addresses; the tensor-tower
sign laws above.  Each with its Lean footprint or its citation, at the point of
use.

\emph{Measured.}  Machine-precision landings across the BSD ladder (rank,
regulator, and the Tate--Shafarevich order as independent jet coordinates,
the obstruction arriving as an exact square integer); the clock-cell phase laws;
the census campaigns on Weil-type classes --- pre-registered statistics,
published nulls included.

\emph{Open, named exactly.}  Two residuals carry the remaining weight, and the
program prices them without cosmetics.  On the analytic side: the
\emph{transfer} --- that every zero of the readout's continuation is registered
by a carrier source --- reduced, losslessly, to one statement about the native
ledger's mean; the reduction is machine-checked, and so is the fact that the
current corpus cannot close it without a genuinely new ingredient.  On the
cycle side: \emph{recognition above grade one} --- from fired coordinate to
constructed cycle --- proven equivalent, modulo the compiled detector, to the
algebraicity statement itself.  Neither is hidden behind the word ``technical.''

\section{Frequently raised objections}

\paragraph{``Isn't this just harmonic analysis with extra words?''}
Harmonic analysis estimates; this bookkeeping \emph{balances}.  The working
outputs are exact integer ledgers, lattice certificates, and closed forms --- not
bounds.  When a residual is $10^{-12}$, the claim is not ``small'' but ``zero,''
and the reason is a theorem about a lattice, not an optimistic constant.

\paragraph{``Is the third dimension real, or a mnemonic?''}
The projection-with-ledger bijection is a theorem, so the geometric object is
exactly as real as the classical one --- they are the same information in two
charts.  The operational answer is better: the geometry \emph{generates} proofs
(the radial squeeze, the clock-chart algebra, the diamond closed form) that were
not being generated from the flat chart, and it forbids mistakes (chart
artifacts) that the flat chart invites.

\paragraph{``Do you assume the big conjectures?''}
No.  Classical inputs enter as named, cited hypotheses at the point of use ---
the same way Gross--Zagier or Jacquet--Shalika enter any paper --- and the two
statements of conjecture strength that remain are isolated as the two open
residuals above, stated with their falsifiers.  The falsification register is
part of the method: pre-committed failure modes, published hits (there have been
retractions, and they stay in the record), and instruments deliberately broken
with wrong clocks to prove they can fail.

\paragraph{``Why sixth roots of unity?''}
Because the carrier is harmonically scaled and $\mathbb Z[\zeta_6]$ is the
hexagonal lattice: the unique scale at which the program's cancellation laws
close exactly, with a rigidity threshold at norm one.  Unit scale is the chart's
convention; harmonic scale is the object's.  Instruments scanned at unit scale
miss real structure --- measured, twice, before it became law.

\section{What kind of thing is this?  (A field, and how you can tell)}

A distinction worth making plainly, because the era invites confusion.  Much of
the strongest machine-assisted mathematics of the moment is \emph{search}: take
the existing method space --- the known transforms, the known potentials, the
known obstructions --- and drive it deeper into known territory than human
patience allowed.  The results are real and sometimes spectacular; the ontology
is inherited.  Old map, mined to exhaustion.

This program is the other thing: it changes what the objects \emph{are}, and
lets the theorems become bookkeeping.  New coordinates --- the carrier, the
clock, the ledgers, the harmonic cells --- are the actual product; theorems are
what the coordinates emit.  The test of such a claim is not any single result
but a set of symptoms, and we state them so the reader can check for them:

\begin{enumerate}
\item \textbf{The questions come from inside.}  A frame is alive when it
generates its own questions faster than it imports them.  The cell-phase laws of
\S4.2 were not asked by anyone; the clock chart produced them the moment
integration moved to carrier time.  The Fermat-grade sign law of \S5 was on no
problem list; the Hodge diamond emitted it.
\item \textbf{The bricks consume each other.}  A bag of tricks does not
compose; a field does.  In the program's recent history, a lattice
classification proven in a morning was load-bearing for a
temperedness--freeze junction by afternoon, which fed a closure theorem, which
fed the sign law.  Results as inputs to results, at a rate outpacing external
input.
\item \textbf{Failure is metabolized, not survived.}  A cancellation
hypothesis tested \emph{anti}-null became a design law about controls.  A
mechanism prediction died cleanly at the data and the corrected mechanism was
sharper.  A vocabulary collision with the classical literature nearly killed a
result and instead produced a better theorem than the one first sought.  The
falsification machinery is part of the engine, not the brakes --- and the
failures stay in the record.
\item \textbf{Someone who is not us runs the machine.}  The decisive symptom,
and the one this document exists to test.  Self-sustainment is complete when an
outsider --- human or machine --- picks up the recipes of \S4 and obtains
something we did not: a derivation we have not written, an application we have
not tried.  We register this as a falsifier, in print: if competent readers,
given these tools, consistently obtain nothing, that is evidence \emph{against}
the claim that this is a field, and it will be reported as such.
\end{enumerate}

\noindent Reader: the fourth symptom is you.  A concrete starter exercise with a
checkable answer: \emph{find the first five zeros of the Riemann zeta function,
as carrier heights}.  Build the eta fiber on the harmonized carrier --- one
phasor per site $n$ at position $x_n=(\pi/3)n$, magnitude $w(n/Z)\,x_n^{-1/2}$
with a $C^\infty$ growth window ($w(0)=1$, $w(1)=0$), phase $-y\log x_n$ so that
every phasor spins as the head climbs --- split into the conductor's two lanes
($n$ odd, $u=+1$; $n$ even, $u=-1$).  Let the head climb and watch
\[
c(Z)\;=\;\frac{|R_{\mathrm{odd}}-R_{\mathrm{even}}|}
              {\max(|R_{\mathrm{odd}}|,|R_{\mathrm{even}}|)}
\]
collapse.  Three things the chart will not have prepared you for.  The zeros are
\emph{heights}: the fifth sits at $Z=e^{32.935\ldots}=2.01\times10^{14}$, and the
famous ordinate is a logarithm because the readout is $y=\log z$.  The
cancellation is balance, not smallness: at that fifth zero both lanes have
modulus $3884.68$ and agree to thirty digits.  And you need not add $10^{14}$
arrows --- above any split $M$ the phase step between neighbouring sites is
$y/n\ll1$, the sites stop resolving, and the tail closes in the Bernoulli ledger
whose first entry is the $B_1=-\tfrac12$ of \S6.2; in the lane \emph{difference}
the continuum part cancels outright, so the whole tail reaches the answer
through the ledger alone.  Five thousand terms, twenty correct digits.  Your
instrument is calibrated when its residual at the located zero is
$|\eta(s-2)|/Z^{2}$ --- the growth window's second moment, reading the readout
two units to the left; that identity is the free check, and it fails loudly if
your window, your lanes, or your spin law is wrong.  Then tell us something we
do not know: put a Dirichlet $L$-function on a second rail, each bank on its own
clock, and read the alignment ledger between the two sets of closures.

\section{How to read the program's papers}

Start anywhere; every paper declares its registers.  When you meet a claim,
ask the three questions the program asks itself: \emph{which chart is this
in?} (carrier-native, or a readout, or an artifact); \emph{what is the
ledger?} (what exact bookkeeping identity is doing the work); \emph{what would
falsify it?} (pre-registered, in print).  If the answer to any of the three is
missing, that is a bug --- report it, and expect the report to be published.

\bigskip
\noindent\emph{This exposition accompanies the main development (the universal
paper and its companions), where every statement above appears at full strength
with its Lean pointer or citation.}
\end{document}
